\section{Match every claim to the evidence}
\label{sec:integrity}

Technical editing begins with preservation. Keep every theorem condition and
reported value intact. Preserve the notation, citations, and stated
uncertainty. Do not make a sentence cleaner by deleting the fact that limits
it.

All examples in this section are schematic. Their theorem conditions, task
names, and reported values are invented to demonstrate an editing decision.

\subsection{Check the source before strengthening prose}

Read the manuscript passage and the evidence it cites before editing a claim.
Check a figure together with its caption, then consult the appendix for omitted
conditions. A result that appears broad in the introduction may depend on a
narrow metric or a specific initialization. A theorem may use a stronger
assumption than the surrounding discussion.

Never invent evidence or imply that planned work is complete. Preserve every
theorem condition and quantitative result. Mark a missing citation, fact, or
control for the authors. Words cannot replace the missing evidence.

\begin{quote}
\textbf{Too broad:} Theorem~3 shows that action chunking prevents compounding
error.

\textbf{Faithful:} Under open-loop stability and bounded disturbances,
Theorem~3 bounds rollout error for horizons up to $H$.
\end{quote}

\subsection{Trace each promise to evidence}

Each abstract claim should appear in the main paper. Contributions in the
introduction should point to visible evidence. A figure caption should make
the same claim as the paragraph that interprets it, and the conclusion should
follow from results already presented.

Run these checks in both directions. A promised contribution with no evidence
is an overclaim. A decisive result that never appears in the abstract or
introduction may mean the paper's contribution hierarchy is wrong. Repair the
hierarchy before polishing sentences.

\begin{quote}
\textbf{Unsupported promise:} Method A consistently outperforms prior methods
across diverse environments.

\textbf{Promise matched to evidence:} At 100{,}000 interactions, Method A
improves mean return over Method B on two of three simulated locomotion tasks
and matches it on the third (Figure~2).
\end{quote}

\subsection{State comparison rules}

State a comparison well enough that a reader can reconstruct its logic. Name
the comparator, metric, and task. Then record any difference in budget,
initialization, or data source. Mechanism claims require tighter controls than
broad performance comparisons because they aim to isolate a cause.

Report failures and negative results with the same precision as successes.
They can identify the regime boundary, reveal sensitivity, or reject an
explanation. Do not hide them in vague limitation language.

\begin{quote}
\textbf{Incomplete:} Method A improves success by eight percentage points.

\textbf{Reconstructable:} Compared with Method B on Task Q, Method A raises
mean success from 68\% to 76\% over 20 evaluation seeds. Both methods use
100{,}000 environment interactions and the same initialization distribution.
\end{quote}

\subsection{Treat limitations as part of the result}

A limitation is useful when it changes how the result should be applied or
interpreted. Name the condition that creates the limit and say what remains
supported. If several limits matter, explain each one separately.

Do not use a broad limitations paragraph to excuse a broad claim. Narrow the
claim where it first appears. Readers should not need to reach the final page
to learn that the central result holds only in one setting.

\begin{quote}
\textbf{Limit buried later:} The method is robust to distribution shift. The
limitations section later reveals that the tests changed only image texture
in simulation.

\textbf{Limit in the claim:} The method maintains its success rate under the
simulated texture shifts tested here. Robustness to dynamics shifts and
real-world observations remains unknown.
\end{quote}
